\section{The complete marker sum}\label{sec:marker-complete}
Here the binary exterior includes every even orbit, including the critical
orbits and their already assembled attachments. Thus no terminal function
is left to be multiplied into the counting comparisons below.

For $a\geq1$ define the full-coordinate ternary fibre
\[
 L_3(a)=\sum_{\substack{K\leq\F_3^a\\
                 \operatorname{pr}_iK=\F_3\ (1\leq i\leq a)}}
                3^{a-\dim K},\qquad L_3(0)=1.
\]
For $g\geq q\geq1$ put
\[
 H_{g,q}=\sum_{\substack{a_1+\cdots+a_q=g\\a_i\geq1}}
                   \prod_{i=1}^q\frac{L_3(a_i)}{a_i!},
 \qquad H_{0,0}=1,
\]
and put $H_{g,0}=0$ for $g>0$.

\begin{proposition}[Exact collapse of repeated signs]
\label{prop:complete-marker-collapse}
Let $\mathcal Y_m$ be a collection of fixed-point-free binary
subgroups on $m$ points, closed under conjugation by $S_m$.
Let $t_q(Y)$ count the $q$-element sets of
actual pair orbits of $Y$ whose characters on $Y$ are distinct.
For $n=m+3g+f-2q$, the number of subgroups obtained by replacing these
selected pairs by $g$ natural $S_3$ orbits in $q$ sign classes, and
adding $f$ fixed points, is exactly
\begin{equation}\label{eq:complete-marker-collapse}
 \frac{n!}{m!}\frac{2^qH_{g,q}}{6^gf!}
                   \sum_{Y\in\mathcal Y_m}t_q(Y).
\end{equation}
Every group so counted recovers $Y$, the selected pairs, the sign
classes, and all its ternary extension data.
\end{proposition}
\begin{proof}
Write the original group as $X\leq T\times S_3^g$, with $T$ its
complete projection on the other moved orbits. The image $B$ of $X$
in $T\times C_2^g$ is a binary group. Its marker signs are nonzero
characters. Group equal signs into classes of sizes $a_1,\ldots,a_q$.
Deleting repeated signs is injective on $B$; realize each retained
sign on a pair. The resulting group is $Y\leq S_m$ and the newly
introduced pairs give the selected set.

Every $B$-invariant subspace of $C_3^g$ splits across these distinct
sign isotypes by an explicit projection. Regard the signs as scalars
in $\F_3$. For distinct signs $\chi,\psi$, choose $b\in B$ with
$\chi(b)\ne\psi(b)$. The operator
\[
 (\chi(b)-\psi(b))^{-1}\bigl(b-\psi(b)I\bigr)
\]
preserves the subspace, is the identity on the $\chi$-isotype and
vanishes on the $\psi$-isotype. Multiplying one such diagonal operator
for every other distinct sign gives the exact $\chi$-projection.
These projections sum to the identity, with each distinct sign counted
once, regardless of its coordinate multiplicity.
In an isotype of dimension $a_i$,
the intersection $K_i$ must have full coordinate projections. For
each such intersection, the lifts are the cocycles into
$\F_3^{a_i}/K_i$. This is a nontrivial sign module, its invariants
are zero, and $H^1(B,\F_3^{a_i}/K_i)=0$. The cocycle set therefore
has exactly $3^{a_i-\dim K_i}$ elements. Hence the literal lift count
is $\prod_iL_3(a_i)$.

The original marker weight is $1/(6^gg!f!)$. Assigning the displayed
marker coordinates to the $q$ selected pairs contributes
$g!/\prod_i a_i!$. A selected pair already has action weight $1/2$,
so replacing the $q$ pairs contributes $2^q$. These give the
coefficient in \eqref{eq:complete-marker-collapse}; the change of
labelled support gives $n!/m!$; here closure under relabelling is
essential. An ordering of the selected pairs only
indexes the summation variables $a_i$; it creates no $q!$ factor.
Projection recovers all the displayed data, proving equality.
\end{proof}

\begin{theorem}[Complete noncanonical marker row]
\label{thm:complete-marker-row}
Write $n=2N+\eps$. For even $m=2M<n$, define
\begin{equation}\label{eq:complete-marker-kernel}
 K_{\rm mark}(n,m)=\frac{c_MG_M}{p_nG_N}
 \sum_{g,f,q}\frac{2^qH_{g,q}}{6^gf!}\binom Mq,
\end{equation}
where the sum is over nonnegative integers satisfying
\[
 n=m+3g+f-2q,\quad
 \begin{cases}q=0,&g=0,\\1\leq q\leq\min(g,M),&g>0,
 \end{cases}\qquad
 D=\frac{g+f-\eps}{2}\geq1.
\]
For all sufficiently large $N$, at both parities,
\begin{equation}\label{eq:complete-marker-row}
 \sum_{\substack{m<n\\m\ \mathrm{even}}}K_{\rm mark}(n,m)
                  \leq2^{-N/10}.
\end{equation}
For collections $\mathcal Y_m$ closed under conjugation by $S_m$,
the complete normalized count of the corresponding marker groups
whose collapsed group lies in $\mathcal Y_m$ is at most
$\sum_mK_{\rm mark}(n,m)|\mathcal Y_m|/L_m$.
\end{theorem}
\begin{proof}
The last assertion follows from
Proposition~\ref{prop:complete-marker-collapse} and
$t_q(Y)\leq\binom Mq$. In particular it does not assume boundedness
of any normalized subgroup count.

Put $\lambda=\log_2 3$ and $k=g-q$. Since $L_3(1)=1$, while
the ternary Gaussian product constant is less than two, for $a\geq2$
\[
 L_3(a)\leq2a\,3^{(a+1)^2/4}
       \leq2^{\lambda(a-1)^2/4+6(a-1)}.
\]
Consequently
\[
 H_{g,q}\leq2^{\lambda k^2/4+6k+g}.
\]
Here $\sum_i(a_i-1)^2\leq k^2$, and the number of compositions of
$g$ is at most $2^g$; dropping $\prod_i a_i!$ is harmless.
For every term of \eqref{eq:complete-marker-kernel},
\[
 M=N-D-k,\qquad 0\leq k\leq\min(2D,N-D),\qquad g\leq2D+1.
\]
The uniform Gaussian estimates leave quadratic exponent
\[
 Q=-\frac{N(D+k)}2+\frac{(D+k)^2}4+\frac{\lambda k^2}4.
\]
We claim $Q\leq-13ND/60$. Set $x=D/N$, $y=k/N$.
The expression $Q/N^2$ is convex in $y$, so it suffices to check
$y=0$ and $y=\min(2x,1-x)$. At $y=0$ its ratio to $x$ is
at most $-1/4$. If $x\leq1/3$, the other endpoint has ratio
$-3/2+(9+4\lambda)x/4\leq-(9-4\lambda)/12$.
If $x\geq1/3$, that ratio is
$(\lambda-1)/(4x)-\lambda/2+\lambda x/4$, a convex function of
$x$ whose endpoint values are $-(9-4\lambda)/12$ and $-1/4$.
The bound follows from $\lambda<8/5$.

The coefficient recurrence gives $c_M/p_n\leq(2N)^{D+k}$.
Also $\binom Mq\leq N^q$, $q\leq3D$, and $6k+g\leq15D$.
The omitted factor $2^q/(6^gf!)$ is at most one. For fixed $D$
there are at most $(2D+2)^2$ choices of $g,f,q$. The complete
fixed-$D$ logarithm is therefore at most
$-13ND/60+O(D\log(N+2))$, with one absolute implied constant.
Eventually it is at most $-ND/6$. The geometric sum over
$1\leq D\leq N$ is at most $2^{-N/10}$.
\end{proof}

\begin{corollary}[Pure markers]\label{cor:pure-marker-sum}
The normalized mass of all subgroups whose even orbits are critical
and whose marker profile has $D\geq1$ is $O(2^{-N/10})$.
This sum includes all actual critical attachments and both parities.
\end{corollary}
\begin{proof}
Collapse yields only critical binary orbits. Theorem~\ref{thm:critical-intro}
bounds their complete normalized count uniformly, including all finite
initial degrees. Apply Theorem~\ref{thm:complete-marker-row} to this
conjugation-invariant family. This argument counts already assembled groups, so it also
includes non-Frattini central attachments.
\end{proof}

\begin{corollary}[One typed binary-error transfer]
\label{cor:complete-marker-error}
Suppose a family consists of fixed points, natural $S_3$ orbits, and
binary orbits, at least one of which is noncritical. Its entire marker
contribution is bounded by
\[
 \sum_{\substack{m<n\\m\ \mathrm{even}}}H_E(n,m)E_{m/2},
\qquad
 H_E(n,m)=K_{\rm mark}(n,m)+\mathbf1_{n=2N+1}
       \bigl(u_N\mathbf1_{m=n-1}+v_N\mathbf1_{m=n-3}\bigr),
\]
with $u_N,v_N$ from \eqref{eq:uv}. There is no scalar remainder.
\end{corollary}
\begin{proof}
The noncritical binary orbit is unchanged by marker collapse, so the
noncanonical part targets the complete conjugation-invariant family
defining $E$. Any additional restrictions only decrease the original count.
The only $D=0$ marker states are no marker in even degree, one fixed
point, and one $S_3$ in odd degree. The even state belongs to the
separately counted complete binary family. The exact odd contributions
are $\alpha_NE_N$ and $\beta_NpE_N$ by
\eqref{eq:complete-marker-collapse}, and
Proposition~\ref{prop:binary-moment} gives the asserted two targets.
Neither four-pair fusion nor duplicate collapse is required to preserve
any earlier exclusion; they preserve a noncritical actual orbit, which
is the defining condition of the complete target $E$.
\end{proof}
